% Generated draft: requires mathematical and source review.
\documentclass[11pt]{article}
\usepackage[margin=25mm]{geometry}
\usepackage{fontspec}
\setmainfont{Latin Modern Roman}
\usepackage{amsmath,amssymb,mathtools,unicode-math}
\setmathfont{Latin Modern Math}
\usepackage{enumitem,needspace,xurl,hyperref}
\setlength{\parindent}{0pt}
\setlength{\parskip}{5pt}
\setlength{\emergencystretch}{3em}
\newcommand{\N}{\mathbb{N}}
\newcommand{\Z}{\mathbb{Z}}
\newcommand{\Q}{\mathbb{Q}}
\newcommand{\R}{\mathbb{R}}
\newcommand{\C}{\mathbb{C}}
\newcommand{\F}{\mathbb{F}}
\newcommand{\A}{\mathbb{A}}
\newcommand{\PP}{\mathbb{P}}
\newcommand{\E}{\mathbb{E}}
\newcommand{\eps}{\varepsilon}
\newcommand{\dd}{\,\mathrm{d}}
\begin{document}
\section*{UnsolvedMath ID: 30002346}
\textbf{Source problem code:} catalogue entry
\subsection{Definitions and mathematical statement}
Let \(\Gamma\) be a finite group acting faithfully and linearly on an ambient space \(V\), and write \(X=V/\Gamma\) for the quotient appearing in the supplied record. Faithfulness means that only the identity element acts trivially on all of \(V\). The base field or base ring, the dimension of \(V\), and the intended construction of the quotient are not specified. Let \(o\in X\) denote the image of the zero vector.

Suppose that \(\pi:Y\to X\) is a crepant resolution. In the usual setting where canonical divisors and their pullback are defined, this means that \(\pi\) is a proper birational resolution with \(Y\) nonsingular and \(K_Y=\pi^*K_X\). The source does not specify the hypotheses needed for this interpretation. Write \(E=\pi^{-1}(o)\) for the fiber over the origin, called the exceptional fiber in the record. If \(E\) is defined over a field \(\kappa\), then \(E(\kappa)\) denotes its set of \(\kappa\)-rational points and \(\#E(\kappa)\) its cardinality.

The proposed universal identity, for quotient data and crepant resolutions satisfying the intended but incompletely specified hypotheses, is
\[
\#E(\kappa)=M(K,\Gamma,-w).
\]
Here \(K\), its relationship to \(\kappa\), the weight \(w\), and the mass expression \(M(K,\Gamma,-w)\) are not defined in the supplied record. In particular, its summation domain, normalization, and finiteness or convergence conditions are absent. The characteristic and ramification conditions intended by the word ``wild'' are also unspecified. Thus the displayed equality preserves the supplied question, but a fully quantified mathematical conjecture requires these missing definitions and hypotheses. The background additionally describes a general motivic equality, without supplying a motivic formulation or its relationship to this point-count identity.
\subsection{Short English statement}
For a quotient by a faithful finite linear group action that admits a crepant resolution, the question asks whether the number of rational points in the fiber over the origin equals the indicated wild McKay mass. The fields, weight, mass definition, and precise admissible quotient data remain unspecified. The supplied status is partially solved; this is a statement-only draft requiring human mathematical review.
\subsection{Sources}
\begin{itemize}
\item Local UnsolvedMath registry, supplied complete source record: permanent ID \(30002346\), exact source problem code \texttt{catalogue entry}, title ``Wild McKay Identity for Crepant Resolutions.'' The collection is ``Oberwolfach Reports Open Problems.'' No individual proposer is identified; the supplied proposed year is \(2013\).
\item Supplied report citation: \emph{Complex Algebraic Geometry (2014)}, Oberwolfach Reports, pp.\ 1563--1627, DOI \texttt{10.4171/owr/2013/27}. \url{https://doi.org/10.4171/owr/2013/27}. The report is supplied as the provenance for a workshop open-problem discussion. Its contents and the attribution of this particular formulation have not been independently verified here. The citation's year \(2014\) and the \(2013\) appearing in the DOI and proposed-year field require reconciliation.
\item Yasuda, \emph{The p-cyclic McKay correspondence via motivic integration}. \url{https://arxiv.org/abs/1208.0132}.
\item Wood--Yasuda, \emph{Mass formulas and quotient singularities}. \url{https://doi.org/10.1093/imrn/rnv074}.
\item Yasuda, \emph{Open problems in the wild McKay correspondence}. \url{https://arxiv.org/abs/2107.07073}.
\item Literature assessment checked at \texttt{2026-08-21}: this is supplied dated curation, not an independently verified assessment. It labels the problem partially solved and reports proofs in important \(p\)-cyclic and symmetric-group or Hilbert-scheme cases, while describing the general motivic equality for arbitrary wild quotient data as conjectural. The scope and hypotheses of those reported results, their relationship to the displayed identity, and time-sensitive status require independent verification.
\item The supplied record labels the statement \texttt{corrected\_verified} and reports a manual semantic, typographic, and LaTeX review claiming a self-contained, meaning-preserving cleanup. These are supplied editorial assertions; they do not resolve the missing mathematical definitions listed above.
\end{itemize}
\end{document}
